\section{The finite scalar verification}\label{sec:computation}

The geometric argument requires the two inequalities
\eqref{cal:scalar-bounds} for the exact functions $G$ and $g_d$.
We now specify a finite calculation that implies them.  Its inputs
are the radial data \eqref{cal:radial-data}, the formulas
\eqref{cal:f0-e}, and the constants defined by the integrals
\eqref{cal:constants}.  Its outputs are interval bounds for these
constants, two strict bounds on a compact rectangle, and a positive
Hessian bound near the exact stationary point.  The estimates of
Section~\ref{tail:section} connect the finite radial calculation to
the infinite series.

Every terminating decimal in this section denotes an exact rational
number.  An interval enclosure always contains the exact mathematical
quantity in question.  In particular, the constants $\lambda,\mu,K$
are not replaced by rounded numerical choices.  Their enclosures are
used in all subsequent operations, while
Lemma~\ref{cal:stationarity} supplies stationarity for their exact
values.

\subsection{Outward arithmetic and derivative intervals}
\label{comp:arithmetic}

The calculation follows the interval inclusion principle: every
operation encloses its exact result, and analytic remainders are added
explicitly. Revol and Rouillier~\cite{RevolRouillier2005} discuss this
principle and the dependence of enclosure width on expressions and
subdivision. We specify the integer operations here, so the proof's
arithmetic assumptions and rounding rules are explicit.

Put $q=2^{144}$.  A scalar interval is represented by integers
$u\le v$, meaning $[u/q,v/q]$.  Rational inputs are rounded down at
the lower endpoint and up at the upper endpoint.  Addition and
negation are exact at this denominator.  Multiplication takes the
least and greatest of the four endpoint products, then divides by
$q$ with outward rounding.  If $0<u\le v$, inversion gives
\[
 \left[\frac{\lfloor q^2/v\rfloor}{q},
       \frac{\lceil q^2/u\rceil}{q}\right];
\]
negative intervals are treated by negation, and division by an
interval containing zero is forbidden.  Integer powers use repeated
multiplication.  For a positive interval, integer square roots give
the enclosure
\[
 \left[\frac{\lfloor\sqrt{uq}\rfloor}{q},
       \frac{\lfloor\sqrt{vq}\rfloor+1}{q}\right].
\]
The extra unit in the upper endpoint is harmless when the upper
square root is an integer.

For logarithms, first extract an integer power of two and reduce the
argument to $[1/4,1]$.  The arguments used below are positive table
intervals or enclosures of fixed positive constants; their widths
allow this reduction.  With $z=(v-1)/(v+1)$, evaluate
\begin{equation}\label{comp:logarithm}
 \log v=2\sum_{j=0}^{199}\frac{z^{2j+1}}{2j+1}+\mathcal R,
 \qquad |\mathcal R|
 \le\frac{2(3/5)^{401}}{401(1-(3/5)^2)}<q^{-1}.
\end{equation}
The same series at $z=1/3$ encloses $\log2$, used to restore the
extracted power.  All summation errors are already enclosed by
outward arithmetic; the displayed remainder is added separately.
Similarly, use
\begin{equation}\label{comp:pi}
 \pi=16\arctan(1/5)-4\arctan(1/239),\qquad
 \arctan(1/v)=\sum_{j=0}^{149}
             \frac{(-1)^jv^{-2j-1}}{2j+1}+\mathcal R_v.
\end{equation}
The alternating-series bound gives
$|\mathcal R_v|\le v^{-301}/301<q^{-1}$ for both values of $v$.
Adding $[-2q^{-1},2q^{-1}]$ to each arctangent enclosure before the
linear combination is sufficient.  These procedures also enclose
$p,R_0,x_0,l$ using their exact defining formulas.

To bound derivatives we use interval jets
\[
 (F_0,F_1,F_2,F_3),\qquad F_j\supseteq F^{(j)}(t)/j!.
\]
Addition, multiplication, inversion, logarithms, and square roots
are performed as formal power series through degree three.
For example, the product coefficient is
$\sum_{r=0}^jF_rG_{j-r}$, and the inverse coefficients follow
recursively from the product being one.  A variable on a radial
interval $T$ is represented by $(T,1,0,0)$; the resulting coefficients
enclose its derivatives at every point of $T$ on the chosen spline
piece.  At a single point, the leading interval encloses that point
and the same derivative entries are used.  After differentiation
only the coefficients supported by the available derivative orders
are used.  The calculation below never requires a fourth derivative
of a spline: third derivatives of $B,D_i$ suffice for the second
derivatives of $f_0,e$ and for the radial quadrature errors.

For a differential equation $Y'=F(t,Y)$ with an enclosed value at
the evaluation point,
its jet coefficients are reconstructed in increasing order by
\begin{equation}\label{comp:jet-ode}
 Y_{j+1}=\frac{[F(t,Y)]_j}{j+1},\qquad 0\le j\le2.
\end{equation}
Thus the enclosed value and the known differential equation supply
the derivatives without numerical differentiation.

Spline polynomials are evaluated by interval Horner arithmetic.
One may first translate a polynomial $P(t)=\sum_{\ell=0}^Na_\ell
t^\ell$ to the midpoint $m$ of its leading interval:
\begin{equation}\label{comp:polynomial-shift}
 P(m+h)=\sum_{j=0}^N
   \left(\sum_{\ell=j}^N\binom{\ell}{j}a_\ell m^{\ell-j}\right)h^j.
\end{equation}
This is an exact polynomial identity, also for jets.  It reduces
interval widths without changing the spline or introducing an
approximation; direct Horner evaluation remains a valid enclosure.

\subsection{Coefficients and the radial table}
\label{comp:radial-table}

Retain the auxiliary mode $0$ and modes $1\le m\le38$.
The two modes with exponent six remain separate in sums involving
$D_i$.  Enclose each $c_m$ by summing over
$|n|,|r|\le D$, where $D=256$ for $m=1$ and $D=48$ otherwise,
and then adding the symmetric error in
\eqref{tail:coefficient-square}, multiplied by $a^{-k_m}/k_m$.
The summands can be evaluated using only integers before that last
factor: if $k=k_m$, $A=2n+r$, and $N=n^2+nr+r^2$, then
\begin{equation}\label{comp:coefficient-summand}
 \operatorname{Re}(n+r/2+i\sqrt3r/2)^{-k}
 =\frac{\displaystyle\sum_{j=0}^{k/2}(-1)^j
       \binom{k}{2j}A^{k-2j}(3r^2)^j}{2^kN^k}.
\end{equation}
The pair $(0,0)$ is omitted.  The resulting enclosure also tests
$|c_1|<1$.

The radial breakpoints are
\begin{equation}\label{comp:breakpoints}
 b,\quad p-.06,\quad p,\quad R_0,\quad R_0+.001,
       \quad R_0+.075,\quad l=R_0+.076.
\end{equation}
Each intervening interval is subdivided by affine interpolation
between its exact endpoints.  For the first interval use
$n=\lfloor8000(u-v)\rfloor+1$ subintervals, and for all later
intervals use $n=\lfloor24000(u-v)\rfloor+1$, where $v,u$ are its
endpoints.  In execution an outward upper bound on the product is
used to choose $n$, which can only add subdivisions.  All true node
positions are enclosed by the corresponding interval arithmetic.
The length of every actual subinterval is bounded above in the
calculation; these upper bounds are used in all error estimates.

For the fixed table, write $[u_i/q,v_i/q]$ for the valid outward
enclosure of its $i$th node, $0\le i\le m$.  The interval
constructor enforces $u_i\le v_i$, and the node arithmetic preserves
this invariant.  After appending the last node and before propagating
any integral state, the program requires
\begin{equation}\label{comp:node-separation}
                    v_i<u_{i+1}\qquad(0\le i<m);
\end{equation}
it throws an exception if a comparison fails.  Thus, on a completed
table construction,
$u_i\le v_i<u_{i+1}\le v_{i+1}$.  Both endpoint arrays are strictly
increasing, and the stored step interval is
\[
 H_i=[(u_{i+1}-v_i)/q,(v_{i+1}-u_i)/q]
\]
with positive lower endpoint.  This also justifies every division by
$H_i$ below.  These facts concern the initialized table and its
unchanged aligned node, point, and cell arrays in the fixed calculation.

On $0<t\le p-.06$, the truncated functions have the monomial form
\[
 \widetilde B(t)=-\log t+\pi t^2/2
       +\sum_{m=1}^{38}a_mt^{k_m},\qquad
 \widetilde B_1(t)=\pi t+\sum_{m=1}^{38}k_ma_mt^{k_m-1},
 \qquad a_m=(-1)^mc_m.
\]
For compact formulas introduce the finite list
\[
 (e_0,v_0)=(2,\pi),\qquad
 (e_m,v_m)=(k_m,k_ma_m)\quad(1\le m\le38).
\]
This list indexes monomials of $\widetilde B_1$, rather than the
sector modes; its exponent $e_0=2$ is unrelated to $k_0=6$.
The starting integrals, including their derivatives, are obtained
from the exact identities
\begin{equation}\label{comp:initial-integrals}
 \begin{split}
 \widetilde S_1(t)&=\frac{t^2}{2}\left(\frac12-\log t\right)
            +\frac\pi8t^4+
                  \sum_{m=1}^{38}\frac{a_mt^{k_m+2}}{k_m+2},\\
 \widetilde S_2(t)&=-2\sum_{i=0}^{38}\frac{v_it^{e_i}}{e_i}
            +\sum_{i,j=0}^{38}\frac{v_iv_jt^{e_i+e_j}}{e_i+e_j},\\
 \widetilde H_k(t)&=\sum_{i=0}^{38}\frac{v_it^{e_i}}{e_i+k}.
 \end{split}
\end{equation}
In particular, the table starts propagation at the exact node
$p-.06$ with these values; preceding nodes and cells use
\eqref{comp:initial-integrals} directly.

There is a useful exact simplification for the sector functions on
this first interval.  Since $b_m(t)=0$, $D_m(t)=-a_mt^{k_m}$ and
$D_0(t)=0$, substitution in \eqref{cal:trial-potential} cancels every
nonconstant mode and gives $H(r,\varphi)=\gamma(r)$ for $r\le t$.
Equivalently, every $Q_i$ vanishes.  Direct radial integration gives
\begin{equation}\label{comp:inner-exact}
 f_0(t)=(\pi t^2-\tfrac12)\log t-\frac{3\pi^2t^4}{8},
 \qquad e(t)=0,\qquad 0<t\le p-.06.
\end{equation}
Indeed,
\[
 \int_0^t\gamma(r)r\,\mathrm dr
 =-\frac{t^2}{2}\log t+\frac{t^2}{4}+\frac{\pi t^4}{8},
 \qquad
 \int_0^t{}^{\rm ren}\gamma'(r)^2r\,\mathrm dr
 =\log t-\pi t^2+\frac{\pi^2t^4}{4}.
\]
Combining these with the source and energy factors in the sector
formula proves \eqref{comp:inner-exact}.  The identity holds for
the full infinite series as well as for every finite truncation.
It may therefore replace the mode sums when evaluating $f_0,e$
and their first two derivatives on the first interval, producing
tighter enclosures.  The full finite-mode values
\eqref{comp:initial-integrals} at $p-.06$ are still required to
start the continuation of $S_1,S_2,H_k$.

For a subsequent radial step $[v,u]$, propagate the integrals by
\begin{equation}\label{comp:propagation}
 \begin{split}
 \widetilde S_1(u)&=\widetilde S_1(v)+
                         \int_v^u r\widetilde B(r)\,\mathrm dr,\\
 \widetilde S_2(u)&=\widetilde S_2(v)+
             \int_v^u\bigl(-2\widetilde B_1(r)
                            +r\widetilde B_1(r)^2\bigr)\,\mathrm dr,\\
 \widetilde H_k(u)&=(v/u)^k\widetilde H_k(v)
              +\int_v^u(r/u)^k\widetilde B_1(r)\,\mathrm dr.
 \end{split}
\end{equation}
For each integrand $Y$, the midpoint rule has the rigorous enclosure
\begin{equation}\label{comp:midpoint}
 \int_v^uY(r)\,\mathrm dr
 \in (u-v)Y((v+u)/2)
   +\left[-\frac{(u-v)^3}{24}\sup_{[v,u]}|Y''|,
           \frac{(u-v)^3}{24}\sup_{[v,u]}|Y''|\right].
\end{equation}
Midpoint values use point jets; second-derivative bounds use a jet
whose leading interval contains $[v,u]$.  The formula applies to
weak second derivatives bounded almost everywhere as well.

For values inside a step, use the integral equations rather than
the midpoint approximation.  If $T=[v,u]$ and $h=u-v$, then
\begin{equation}\label{comp:cell-integrals}
 \begin{split}
 \widetilde S_j(T)&\subseteq \widetilde S_j(v)
                                  +[0,h]\,\operatorname{range}_T Y_j,
                         \qquad j=1,2,\\
 \widetilde H_k(T)&\subseteq (v/T)^k\widetilde H_k(v)
                  +[-h\|\widetilde B_1\|_{0,T},
                     h\|\widetilde B_1\|_{0,T}],
 \end{split}
\end{equation}
where $Y_1=r\widetilde B$ and
$Y_2=-2\widetilde B_1+r\widetilde B_1^2$.
For the last inclusion, $(r/t)^k\le1$ whenever $v\le r\le t\le u$.
Recover the positive jet orders from
\begin{equation}\label{comp:radial-odes}
 \widetilde S_1'=t\widetilde B,\qquad
 \widetilde S_2'=-2\widetilde B_1+t\widetilde B_1^2,\qquad
 \widetilde H_k'=\widetilde B_1-k\widetilde H_k/t.
\end{equation}
Substitution into \eqref{cal:f0-e} produces both node enclosures and
cell enclosures of the ordered six-tuple
\begin{equation}\label{comp:table-tuple}
 (\widetilde f_0,\widetilde f_0',\widetilde f_0'',
          \widetilde e,\widetilde e',\widetilde e'').
\end{equation}
No numerical differentiation of table values is used.

Separately, subdivide each interval in
\eqref{comp:breakpoints} into steps of length at most $.00022$.
Direct interval jets for $\widetilde B_1$ and $\widetilde D_i$
test the two finite-core inequalities
\eqref{tail:finite-core-premise}.  Together with $|c_1|<1$, these
are the finite premises of the truncation analysis; their verification
does not use any infinite-mode error bound.  Only after these tests
pass is \eqref{tail:uniform-error} invoked.

The spline data are $C^2$ and have bounded one-sided third
derivatives.  Accordingly $f_0,e$ are $C^1$ with locally bounded weak
second derivatives on $(b,\infty)$.  A radial cell uses its own spline
formula at its endpoints; either one-sided second derivative is
valid for its interior.  When a query meets a knot, the hull includes
both adjacent cells.  Thus no jump of a second derivative is mistaken
for global classical differentiability.

\subsection{Queries for the infinite-mode functions}
\label{comp:lookup}

Let $\epsilon=[-10^{-11},10^{-11}]$, and let $T$ be a valid outward
enclosure of the queried true radii.  In the fixed production calls
every such radius is at least $b$, although $T$ may extend below $b$.
For derivatives, take the componentwise hull of
\eqref{comp:table-tuple} over every radial cell meeting $T$ and add
$\epsilon$ to each component.  By \eqref{tail:uniform-error}, this
encloses the full infinite-mode tuple $(f_0,f_0',f_0'',e,e',e'')$
at every queried true radius in $T\cap[b,l]$.  The exact continuation
below covers queried true radii at least $l$.
Here is the cell-selection argument for the implementation.  In
integer endpoint units write the valid query as $[a,b']$, let $j$ be
the first index with $v_j\ge a$, and let $k$ be the first index with
$u_k>b'$, using $m+1$ if either index is absent.  The closed envelope
of cell $i$ is $[u_i,v_{i+1}]$.  It meets $[a,b']$ exactly when
$v_{i+1}\ge a$ and $u_i\le b'$.  By the ordering proved above, these
conditions are equivalent to $i\ge j-1$ and $i<k$.  Therefore the
binary searches on the upper and lower endpoint arrays select exactly
\[
 \mathrm{im}=\max(0,j-1)\le i<\min(m,k)=\mathrm{ix}.
\]
The non-strict intersection inequalities include both cells at a
shared endpoint.  Every true cell meeting the query lies in its
envelope and is therefore selected; any extra cells only widen the
hull.  This is a statement about the fixed initialized table, not a
general-purpose lookup interface.

For a query requiring values alone, improve the enclosure by linear
interpolation of the endpoint values.  On a radial cell $[v,u]$ put
\[
 X=[0,1]\cap\frac{T-v}{u-v},\qquad
 L_Y(T)=Y(v)+X\bigl(Y(u)-Y(v)\bigr),
 \quad Y=\widetilde f_0,\widetilde e.
\]
An empty intersection is discarded.  The interpolation error is
\begin{equation}\label{comp:radial-interpolation}
 |Y(t)-L_Y(t)|\le\frac{(u-v)^2}{8}
                                   \sup_{[v,u]}|Y''|.
\end{equation}
The cell enclosure supplies the derivative bound.  Add this
symmetric error and $\epsilon$, then take the hull over all relevant
cells.  Shared coefficient uncertainties need not be independent:
interval operations enclose all their possible values and hence
also their correlated exact values.
In particular, the positive stored interval $H_i$ contains the true
step $u-v$, so outward evaluation of the quotient defining $X$
contains the normalized coordinate of every true queried point in
that cell.  Intersecting with $[0,1]$ cannot discard such a point.

Queries that meet $[l,\infty)$ also include the exact continuation
\eqref{tail:exterior-formula}.  Enclose its two constants by the
last node value plus $\epsilon$ and by
\begin{equation}\label{comp:constant-B}
 B(l)\in-\log R_0+\pi R_0^2/2+
                    \sum_{m=1}^{38}V_m(R_0)+\epsilon.
\end{equation}
Thus the continuation contributes the tuple
\begin{equation}\label{comp:tail-tuple}
 \bigl(f_0(l)-\pi B(l)(T^2-l^2),\,
       -2\pi B(l)T,\,-2\pi B(l),\,0,0,0\bigr).
\end{equation}
Using the entire query interval in this expression is conservative
even if it extends below $l$, since the relevant exterior subset is
contained in it.  This rule propagates the constant uncertainties
through the exact formula and asserts no radius-independent error
on the unbounded tail.

\subsection{Enclosing the calibrated constants}
\label{comp:constants}

We evaluate the exact integrals defining $\lambda$ in
\eqref{cal:constants}.  Put $s_0=.01$ and
$t(s)=\sqrt{p^2+s^2}$.  Taylor's formula at $p$ removes the apparent
division by zero in the first integral.  Specifically,
\begin{equation}\label{comp:lambda-origin}
 \begin{split}
 \int_0^{s_0}p\frac{f_0(t(s))-f_0(p)}{s^2}\,\mathrm ds
 \in{}& f_0'(p)\left(\frac{s_0}{2}
           -\frac{s_0^3}{24p^2}
                      +\left[0,\frac{s_0^5}{80p^4}\right]\right)\\
 &+\frac{s_0^3}{6}
   \left[\frac{p}{(\sqrt{p^2+s_0^2}+p)^2},\frac1{4p}\right]
          \operatorname{range}_{[p,t(s_0)]} f_0''.
 \end{split}
\end{equation}
Indeed, $t-p=s^2/(t+p)$, while
\[
 \frac{p}{t+p}=\frac12-\frac{s^2}{8p^2}
                         +\left[0,\frac{s^4}{16p^4}\right].
\]
The quadratic Taylor remainder of $f_0$ supplies the second line of
\eqref{comp:lambda-origin}.  The derivative range comes from the
radial cells with the infinite-mode error added.

The implementation obtains this range by a direct cell scan over
$[p,t(s_0)]$, and obtains the derivative bounds below by a second
direct scan over $[p,R_0]$.  Both use the same closed-cell selection
rule, with outward endpoint bounds, but neither adds the exterior
formula.  This is sufficient here: $p^4=1/12$ gives $p>1/2>b$,
$x_0=R_0/2>s_0$, and
\[
              t(x_0)=\sqrt{p^2+(R_0/2)^2}=R_0<l.
\]
Hence both true scan domains are contained in the finite table
$[b,l]$.  A nonempty scan alone would not justify a query extending
beyond that table.

On $[s_0,x_0]$ integrate
$p(f_0(t(s))-f_0(p))/s^2$ by a midpoint sum with $45000$ equal
steps.  Integrate $e(t(s))/p$ on $[0,x_0]$ by a separate midpoint
sum with $45000$ equal steps.  Obtain integer upper bounds
\[
 C\ge\sup_{[p,R_0]}|f_0'|,\quad
 C'\ge\sup_{[p,R_0]}|f_0''|,\quad
 E\ge\sup_{[p,R_0]}|e'|,\quad
 E'\ge\sup_{[p,R_0]}|e''|
\]
from the cell enclosures plus $\epsilon$.  The two integrands have
second derivatives bounded respectively by
\begin{equation}\label{comp:lambda-derivatives}
 \frac{C'}p+\frac{8C}{s_0^2},\qquad
 \frac{E'}p+\frac E{p^2}.
\end{equation}
For the first bound, set $V=f_0(t(s))-f_0(p)$; then
\[
 |V|\le Cs^2/(2p),\qquad
 V''=f_0''(t)(s/t)^2+f_0'(t)p^2/t^3,
\]
and differentiate $pV/s^2$ twice.  The resulting terms are
$p(V''/s^2-4V'/s^3+6V/s^4)$, bounded by the first expression in
\eqref{comp:lambda-derivatives}.  The second follows by twice
differentiating $e(t(s))/p$.
For an interval of length $L$ and equal step length $h$, add the
symmetric error $Lh^2\sup|Y''|/24$ to each complete midpoint sum.
Together with \eqref{comp:lambda-origin}, the endpoint terms and
the factors in \eqref{cal:constants}, this encloses $\lambda$.
The same formula then encloses $\mu$.

To enclose $K=2\int_0^{x_0}G(p,s)\,\mathrm ds$, choose an integer
$n$ and equal steps of length $h=x_0/n\le.0001$, and use their
midpoint sum.  Independently enclose
$G_{ss}(p,s)$ on a partition with step length at most $.0002$.
If $h$ is the midpoint step length and $L_K$ bounds this second
derivative, the enclosure is
\begin{equation}\label{comp:K-quadrature}
 K\in2\left(\sum_jhG(p,s_j)
                   +[-x_0h^2L_K/24,x_0h^2L_K/24]\right).
\end{equation}
The factor two multiplies both the sum and its error.

The first sign tests are
\begin{equation}\label{comp:tail-signs}
 \begin{gathered}
 -6.731<\lambda<-6.729,\qquad K<0,\\
 f_0(l)-\lambda l^2/2-\mu>0,\qquad
                  -\pi B(l)-\lambda/2>0.
 \end{gathered}
\end{equation}
Each strict inequality is accepted only when the corresponding
interval endpoint proves it.  The bound on $\lambda$ also serves
as a check on the normalization.  The last two signs and
\eqref{tail:exterior-formula} give
\begin{equation}\label{comp:positive-tail}
 f(t)=f(l)+\bigl(-\pi B(l)-\lambda/2\bigr)(t^2-l^2)>0,
 \qquad G(d,s)>0\quad\hbox{if }t=\sqrt{d^2+s^2}\ge l.
\end{equation}
Here $e(t)=0$ and $d/t^2>0$.

\subsection{The compact rectangle and its derivative bounds}
\label{comp:rectangle}

Because $l<.697$ and $\sqrt{b^2+.65^2}>l$, the signs
\eqref{comp:positive-tail} reduce the remaining work to
\begin{equation}\label{comp:domain}
 \mathcal D=[b,.697]\times[0,.65].
\end{equation}
For $d\ge.697$, $G(d,s)>0$ for every $s\ge0$ and $g_d(x)\ge0$.
For $d\ge b$ and $x\ge.65$, the primitive is increasing in $x$.
The negative part of $G$ vanishes for $s\ge.65$.

We obtain derivatives of $G$ by differentiating its defining
formula, using the full radial query enclosures.  The following
identities specify all needed factors.  Put
\[
 t=\sqrt{d^2+s^2},\quad A=d/t^2,\quad
 Z=-s^2/(dt^2)=A-1/d,\quad r_d=d/t,\quad r_s=s/t,
\]
and write $f'=f_0'-\lambda t$, $f''=f_0''-\lambda$.
Then
\begin{equation}\label{comp:rational-derivatives}
 \begin{aligned}
 A_d&=(1-2d^2/t^2)/t^2,&Z_d&=A_d+1/d^2,\\
 A_s&=-2ds/t^4,&Z_s&=A_s,\\
 A_{dd}&=2d(d^2-3s^2)/t^6,&Z_{dd}&=A_{dd}-2/d^3,\\
 A_{ss}&=-2d/t^4+8ds^2/t^6,&Z_{ss}&=A_{ss}.
 \end{aligned}
\end{equation}
For $v=d$ or $v=s$ let $r_v$ have the meaning above.  The derivative
formulas are
\begin{equation}\label{comp:G-derivatives}
 \begin{split}
 G_v={}&A_vf+Af'r_v+Z_ve+Ze'r_v,\\
 G_{vv}={}&A_{vv}f+2A_vf'r_v
       +A\bigl(f''r_v^2+f'(1-r_v^2)/t\bigr)\\
 &+Z_{vv}e+2Z_ve'r_v
       +Z\bigl(e''r_v^2+e'(1-r_v^2)/t\bigr).
 \end{split}
\end{equation}
These formulas hold on each smooth piece and almost everywhere
across the radial knots.  No mixed second derivative of $G$ is
needed.

Partition both sides of \eqref{comp:domain} into intervals of
length at most $.0008$.  Interval substitution in
\eqref{comp:G-derivatives} on every closed rectangle supplies
global upper bounds
\begin{equation}\label{comp:global-derivative-bounds}
 L_1\ge\|G_d\|_{\infty,\mathcal D},\quad
 L_2\ge\|G_s\|_{\infty,\mathcal D},\quad
 L_3\ge\|G_{dd}\|_{\infty,\mathcal D},\quad
 L_4\ge\|G_{ss}\|_{\infty,\mathcal D}.
\end{equation}
The value-only interpolation is not used to obtain these derivative
bounds.

\subsection{The primitive away from its stationary point}
\label{comp:outer-grid}

Let
\begin{equation}\label{comp:inner-domain}
 \mathcal D_0=[.52528,.54930]\times[.29420,.32621].
\end{equation}
Its interior contains $(p,x_0)$.  To verify the primitive bound on
the complement, partition the $d$ interval at the two vertical
sides of $\mathcal D_0$, and partition the $s$ interval at its two
horizontal sides.  Subdivide these pieces into steps with maximal
lengths
\[
 D_h\le.00019,\qquad S_h\le.00018.
\]
Every cell therefore lies either in $\mathcal D_0$ or outside its
interior.  For each $d$ node, form successive midpoint sums from
$0$ to every $x$ node, using the actual step lengths:
\[
 Q(d,x)=\sum_{[s_j,s_{j+1}]\subseteq[0,x]}
          (s_{j+1}-s_j)
                     G\left(d,\frac{s_j+s_{j+1}}2\right).
\]
The sum at $x=0$ is zero and is tested as well.  The value queries
in Section~\ref{comp:lookup} may be used here.  At every node
outside the interior of $\mathcal D_0$, enclose $2Q(d,x)-K$ and
let $m_*$ be a lower bound for all these enclosures.

The quadrature error at grid nodes is bounded by $E_0$, and the
additional interpolation error between nodes by $E_I$, where
\begin{equation}\label{comp:outer-errors}
 E_0=\frac{2(.65)S_h^2L_4}{24},\qquad
 E_I=\frac{2\bigl(D_h^2(.65)L_3+S_h^2L_2\bigr)}8,
\end{equation}
Indeed, the bilinear interpolation estimate follows from
\[
 \partial_d^2(2g_d(x)-K)=2\int_0^xG_{dd}(d,s)\,\mathrm ds,
 \qquad \partial_x^2(2g_d(x)-K)=2G_s(d,x).
\]
Apply the one-dimensional interpolation error in each variable
successively; the resulting bounds add and require no mixed
derivative.  Thus the finite assertion
\begin{equation}\label{comp:outer-assertion}
                         m_*-E_0-E_I>0
\end{equation}
proves $2g_d(x)>K$ throughout
$\mathcal D\setminus\operatorname{int}\mathcal D_0$.
All the derivative estimates are weak derivative estimates, so
the same conclusion holds when a quadrature segment or rectangle
crosses a radial knot.

\subsection{The negative part on the whole rectangle}
\label{comp:negative-part}

At every $d$ node, including those between the vertical sides of
$\mathcal D_0$, use the same $s$ partition to enclose
\[
 N(d)=\sum_j(s_{j+1}-s_j)
       \max\left\{0,-G\left(d,\frac{s_j+s_{j+1}}2\right)\right\}.
\]
Let $N_*$ be an upper bound for all these sums.  The function
$y\mapsto\max\{0,-y\}$ is $1$-Lipschitz.  On a subinterval of
length $h$, its midpoint error is therefore at most $L_2h^2/4$:
integrate $L_2|s-s_{\rm mid}|$.  Summing the errors gives at most
$.65S_hL_2/4$.  Every $d$ lies within $D_h/2$ of a node; changing
$d$ to that node changes the integral by at most
$.65D_hL_1/2$.  Consequently the assertion
\begin{equation}\label{comp:negative-assertion}
 N_*+K+.65\left(\frac{S_hL_2}{4}+\frac{D_hL_1}{2}\right)<0
\end{equation}
proves
\[
 \int_0^{.65}\max\{0,-G(d,s)\}\,\mathrm ds<-K
                   \qquad(b\le d\le.697).
\]
Together with \eqref{comp:positive-tail}, this gives the entire
negative-part inequality in \eqref{cal:scalar-bounds}.

\subsection{Convexity in the stationary rectangle}
\label{comp:hessian}

It remains to prove the primitive inequality on $\mathcal D_0$.
Partition its $d$ interval into steps of length at most $.00004$.
Partition $[0,.32621]$ into steps of length at most $.00008$,
including $.29420$ as a node.  The accumulation in $s$ must begin
at zero even though the Hessian is tested only above $.29420$.

For a $d$ interval $D$ and an $s$ interval $S_j$, use the full
derivative queries to enclose
\[
 F_1(D,S_j)\supseteq G_{dd}(D,S_j),\quad
 F_2(D,S_j)\supseteq G_d(D,S_j),\quad
 F_3(D,S_j)\supseteq G_s(D,S_j).
\]
For a cell $D\times S_j$ meeting $\mathcal D_0$, enclosures of
the three Hessian entries of $(d,x)\mapsto g_d(x)$ are
\begin{equation}\label{comp:hessian-entries}
 \begin{split}
 A_j&=\sum_{i<j}|S_i|F_1(D,S_i)
                             +[0,|S_j|]F_1(D,S_j),\\
 B_j&=F_2(D,S_j),\qquad C_j=F_3(D,S_j).
 \end{split}
\end{equation}
The first line is direct interval integration, not midpoint
quadrature: it bounds both the completed subintervals and the
partial final subinterval.  Test, on every such cell, that
\begin{equation}\label{comp:hessian-assertions}
 \inf A_j>0,\qquad \inf C_j>0,\qquad
                         \inf(A_jC_j-B_j^2)>0.
\end{equation}
Interval multiplication encloses the determinant even when its
factors are correlated.  These inequalities imply positive
definiteness, and in particular positive semidefiniteness, of the
Hessian almost everywhere on $\mathcal D_0$.

The function $(d,x)\mapsto g_d(x)$ belongs to
$W^{2,\infty}$ on a neighbourhood of this rectangle.  To justify
the passage from Hessian bounds to convexity, mollify on each
smaller interior rectangle: the mollified Hessian is positive
semidefinite, so the mollified functions are convex.  Local
uniform convergence proves convexity in the interior, and
continuity includes the boundary.  By
Lemma~\ref{cal:stationarity}, $(p,x_0)$ is an interior stationary
point.  Hence
\begin{equation}\label{comp:inner-conclusion}
          2g_d(x)\ge2g_p(x_0)=K\qquad((d,x)\in\mathcal D_0).
\end{equation}
This step uses the exact stationary point and does not infer
convexity or stationarity from sampled function values.

\subsection{Finite assertions and the scalar conclusion}
\label{comp:conclusion}

For clarity, the finite procedure has the following fixed
parameters and obligations.
\begin{center}
\begin{tabular}{ll}
\hline
Quantity & Parameter or required assertion\\
\hline
Arithmetic & Denominator $2^{144}$, jet order $3$\\
Retained modes & $i=0,1,\ldots,38$\\
Coefficient squares & $D=256$ for $m=1$, $D=48$ otherwise\\
Core mesh & At most $.00022$ on each spline interval\\
Core assertions & $|c_1|<1$ and \eqref{tail:finite-core-premise}\\
Radial table & Densities $8000$ and $24000$ as specified above\\
Radial node ordering & Strict separation \eqref{comp:node-separation}
                       at every adjacent pair\\
$\lambda$ integrals & $45000$ midpoint steps each; $s_0=.01$\\
$K$ integral & Steps at most $.0001$; derivative mesh $.0002$\\
Exterior signs & All assertions in \eqref{comp:tail-signs}\\
Global derivative mesh & At most $.0008$ in both coordinates\\
Outer and negative-part mesh & $D_h\le.00019$, $S_h\le.00018$\\
Outer and negative-part assertions & \eqref{comp:outer-assertion},
                                      \eqref{comp:negative-assertion}\\
Hessian mesh & $d$ steps $.00004$, $s$ steps $.00008$ or smaller\\
Hessian assertions & \eqref{comp:hessian-assertions} on every cell\\
\hline
\end{tabular}
\end{center}

\begin{proposition}[Meaning of the finite assertions]
\label{comp:finite-soundness}
Apply the outward interval procedure specified in this section to
the radial data \eqref{cal:radial-data}, retaining modes $0$ through
$38$, with the parameters in the table.  If every listed assertion
holds, its enclosures establish the finite-core premise
\eqref{tail:finite-core-premise}, the exterior positivity
\eqref{comp:positive-tail}, the primitive bound outside the interior
of $\mathcal D_0$, the negative-part bound for every $d\ge b$, and
the convexity bound \eqref{comp:inner-conclusion}, all for the exact
infinite-mode functions.
\end{proposition}
\begin{proof}
The coefficient sum and its lattice tail enclose each retained
coefficient.  The direct core tests concern only finitely many
modes, so they imply the premise needed for
\eqref{tail:uniform-error}.  Equations
\eqref{comp:initial-integrals}--\eqref{comp:radial-odes} and
the midpoint bound prove the node and cell table enclosures by
induction through every radial step.  The node-separation guard supplies
the positive steps and ordered search arrays used by the table and
query arguments.  The query rules then enclose
the exact infinite-mode functions and their first two derivatives,
including the constant-region continuation.

The origin expansion and midpoint errors prove the enclosures of
$\lambda,\mu,K$.  The signs \eqref{comp:tail-signs} imply the
positive tail.  The derivative bounds, quadrature errors and
bilinear interpolation prove the outer primitive assertion;
the Lipschitz estimate proves the negative-part assertion.
Finally, \eqref{comp:hessian-entries} encloses the complete
Hessian, whose positivity and exact stationarity give
\eqref{comp:inner-conclusion}.  Every stage uses only enclosures
established before that stage.
\end{proof}

\paragraph{Finite bounds.}
The program \texttt{certificate.cpp} implements this procedure using
\texttt{core.hpp} and \texttt{grids.hpp}. All three files and an
execution wrapper are included in the \texttt{verification} directory.
One completed guarded run is documented in
\texttt{verification/README.md}, which binds the code and formula hashes,
run metadata, saved integer pairs, and their interpretation below.
The wrapper compiles and runs the computation; it does not read the TeX
manuscript. The source-to-code and analytic correspondences are separate
arguments.
The interval denominator is $2^{144}$. The table gives outward-rounded
projections, while interval assertions use uncoarsened endpoints.
\begin{center}
\begin{tabular}{lr}
\toprule
Quantity & Recorded projection\\
\midrule
$\lambda$ & $[-6.7300312,-6.7300275]$\\
$\mu$ & $[.4138830,.4138838]$\\
$K$ & $[-.0034823,-.0034809]$\\
$B(l)$ & $[1.0464643,1.0464644]$\\
$f_0(l)$ & $[-1.2135966,-1.2135965]$\\
$f(l)$ & $[.0044755,.0044772]$\\
$-\pi B(l)-\lambda/2$ & $[.0774490,.0774509]$\\
$m_*-E_0-E_I$ & $\ge .0000120$\\
Negative-part left side of \eqref{comp:negative-assertion}
 & $\le -.0004585$\\
$g_{dd}$ on $\mathcal D_0$ & $\ge1.0529222$\\
$g_{xx}$ on $\mathcal D_0$ & $\ge .0128335$\\
$\det\nabla^2g$ on $\mathcal D_0$ & $\ge .0086252$\\
\bottomrule
\end{tabular}
\end{center}
Here the determinant is enclosed on each individual rectangle before
the bounds are combined; it is not computed from separately combined
entry ranges. The global derivative bounds, rounded outwards for
display, are
\[
 L_1\le7.4636170,\qquad L_2\le2.5091315,\qquad
 L_3\le143.7890461,\qquad L_4\le315.1489689.
\]
These derivative numbers are projections of the computed hulls, not
separate executable threshold assertions. For the identified sources,
under the documented runtime assumptions, the completed run verifies the
listed finite assertions using uncoarsened interval endpoints.
Proposition~\ref{comp:finite-soundness} explains why these assertions
imply the required bounds for the exact infinite-mode functions:
the interval formulas and range-coverage arguments identify the
quantities enclosed, while the analytic truncation estimates supply
the omitted modes. Neither the rounded table nor the execution record
alone supplies those mathematical correspondences.

\begin{theorem}[Scalar inequalities from the finite verification]
\label{comp:scalar}
The exact constants and functions of
\eqref{cal:constants} and \eqref{cal:primitive} satisfy
\[
 K<0,\qquad 2g_d(x)\ge K\quad(d\ge b,\ x\ge0),\qquad
 \int_0^\infty\max\{0,-G(d,s)\}\,\mathrm ds\le-K
                                             \quad(d\ge b).
\]
\end{theorem}
\begin{proof}
The documented finite verification and the correspondences just
established supply the hypotheses of
Proposition~\ref{comp:finite-soundness}. The first assertion is
part of \eqref{comp:tail-signs}.
The outer-grid assertion and \eqref{comp:inner-conclusion}
cover all of $\mathcal D$.  For $d\in[b,.697]$ the positive
tail makes the primitive increasing after $x=.65$, so the
primitive bound extends to every $x\ge0$.  For $d\ge.697$,
$G(d,s)>0$ for all $s\ge0$, whence $g_d(x)\ge0>K/2$.
The negative-part estimate follows from
\eqref{comp:negative-assertion} on the rectangle; its integrand
vanishes outside the rectangle by
\eqref{comp:positive-tail}.  These are exactly
\eqref{cal:scalar-bounds}.
\end{proof}
